\section{路径、环与连通性}
\label{sec:04-Discrete-Mathematics/03-Graph-Theory/02}

现在你有一张图，500 个顶点和若干条边。你问：从顶点 A 能走到顶点 B 吗？如果能，最少要经过几条边？

这个问题在实践中出现的频率远超你的预期。路由协议在计算下一跳，社交平台在推荐``二度人脉''，编译器在做数据流分析——它们都在反复追问谁可以到达谁。但你不能对每一对顶点从头开始搜索：一个中等规模的图，顶点对的数量是顶点的平方，暴力枚举等于自杀。

\subsection{把一步邻接累积成可达关系}

先从概念开始清理。邻接是一步就能到，可达是可以经过多步。

游走是顶点和边交替的序列，允许重复访问同一个顶点、重复走同一条边。迹不允许重复边，但允许重复顶点。路径连顶点也不允许重复。环是起终点重合、且内部顶点互不相同也不与起终点重复的闭合路径。

这些区分的用途在于：``最短路一定可以取为路径''（重复顶点意味着绕了圈子，删掉圈子只会更短或等长），而``Euler 回路必须用尽每条边''关心的是迹而非路径。不同教材在用词上偶有出入，每次使用时声明你的定义就不会造成混淆。

无向图上的可达性是一个等价关系——自反（每个顶点到自己有一条长度为 0 的平凡路径）、对称（无向边可以双向穿越）、传递（路径可以拼接）。因此，可达关系把顶点集划分成连通分量：分量内部的任意两点互相可达，不同分量之间没有路径相连。图连通意味着只有一个分量，全体顶点属于同一张网。

这个划分把一个全局问题（``网络是否分裂''）转化成了等价类结构：检查每个顶点属于哪个分量即可。加一条连接两个不同分量的边，分量总数减少 1；在同一分量内部加边，分量数不变。这句话虽然简单，却是一个反复出现的计数工具。

\subsection{最短路拒绝绕圈子}

无权图中，最短路的长度就是经过的边数。关键观察是：任何一条候选路径如果重复了某个顶点，它一定在中间兜了一个闭合的环。把这个环删掉，剩下的仍然是一条从起点到终点的路径，而且边数更少。因此，最短路可以取为简单路径（不重复顶点）。

广度优先搜索（BFS）正是按照这个逻辑逐层推进的：从源点出发，先把所有距离为 1 的邻居标记为``已发现、距离 1''，再扫描这些邻居的未访问邻居标记距离 2，以此类推。当一个顶点第一次被 BFS 发现时，到达它的那条路径就是最短的——因为距离更小的层已经全部扫描完毕，任何尚未发现的顶点都不可能通过更短的路径到达。

加权图的情况更微妙。如果所有边权非负，依然可以用 Dijkstra 算法的贪心策略：每次从尚未确定最短距离的顶点中取出距离估计最小的那个，确认它的距离不会再改善，然后更新它的邻居。这个``确认后锁定''的操作依赖边权非负——一旦所有待处理的顶点都有不小于当前最小值的距离估计，已经锁定的顶点的确不可能被后来的顶点绕路改善。

但当边权可以为负时，这个逻辑碎裂了。你可能需要经过一个远端的负权边``倒贴''一段距离，使一条看似绕远的路径反而更短。此时 Dijkstra 的贪心锁定不再安全，Bellman-Ford 算法的反复松弛才给出正确保证。

\subsection{Euler 盯着边，Hamilton 盯着顶点}

这两个问题常常被放在一起讲，但它们的难度天差地别，原因是它们关心的对象完全不同。

Euler 回路要求经过每条边恰好一次并回到起点。连通无向图存在 Euler 回路的充分必要条件是每个顶点的度数均为偶数。必要性直观：每次进入一个中间顶点，就必须从另一条边离开，入和出成对消耗邻边，要求度数为偶数。充分性可以通过不断拼接闭合回路来构造：从任意顶点出发沿未用过的边走，度数全为偶保证走到无路可走时必然回到了起点，此时得到了一个闭合回路；删去已用边后剩余子图度数仍全为偶，重复这个过程，最后在公共顶点处把所有回路拼接成一条。

允许起点和终点不同时，恰有两个奇度顶点充当端点。条件干净、构造算法明确——Euler 回路是一个完美解决的问题。

Hamilton 回路要求经过每个顶点恰好一次并回到起点。听起来和 Euler 回路对称——Euler 回路遍历边，Hamilton 回路遍历顶点——但在计算上，判定一个图是否有 Hamilton 回路是 NP-完全的。没有像``所有度数都是偶数''这样简洁的充要条件。不要因为两个名字都谈``走遍一张图''就把它们当成同类问题。

\subsection{删点或删边如何破坏连通}

割边是一条边，删掉它之后图的连通分量数会增加。割点是一个顶点，类似的效果。在网络设计中，割边和割点对应单点故障——一个关键服务器宕机就导致内网分裂。

比``是否存在割边''更深入的问题是``两点之间存在多少条内部不相交的路径''。多条不相交的路径意味着冗余——即使一条路径上的边或顶点全部失效，仍有替代路线。Menger 定理精确地建立了这种冗余与最小割的关系：任意两点间内部顶点不相交的路径的最大条数，等于把这两点分开所需删除的最小顶点数。这是连通性从``是或否''走向``有多强韧''的关键升级。

一句话收束：连通只保证``至少有一条路''，不保证路短、路多或路宽。需要这些保证的时候，你得去查最短路径算法、Menger 数和网络流。

下一篇讨论一种极端情况——图刚好连通到不能再删任何一条边，多一条边则出环。这种临界结构叫树。
